\chapter{1990 Nobel Prize in Economics}
\textbf{Laureates: }Harry Markowitz (USA), Merton Miller (USA), William Sharpe (USA)

\section*{Reason for Award}
Portfolio theory derivatives pricing capital structure

\section{What They Did}
Harry Markowitz created modern portfolio theory and mean-variance optimization establishing diversification principle reducing unsystematic risk through combination of assets Markowitz quadratic optimization approach selecting portfolios efficient frontier maximizing return for given minimizes risk level William Sharpe developed Capital Asset Pricing Model CAPM relating expected returns systematic risk beta factor providing simple elegant framework pricing risky assets Sharpe showed market portfolio contains all assets investors hold efficient frontier reduces to straight line representing tradeoff risk return Merton Miller developed MM theorem on capital structure irrelevance showing firm value independent financing decisions debt-equity split perfect markets MM theorem laid foundation for corporate finance revealing financing decisions irrelevant absent imperfections Later extended Miller work showing effects taxes bankruptcy costs agency costs determine optimal capital structure

\section{Impact on Economic Thinking}
Markowitz revolutionized investment management providing systematic framework constructing portfolios balancing risk-return tradeoffs his work established modern investment practice financial advisory asset management mutual funds pension plans Sharpe CAPM became standard tool evaluating investment performance measuring risk adjusting returns Performance evaluation benchmarks derive directly from CAPM Miller MM theorem initiated corporate finance research examining effects financing decisions on firm value Their combined contributions transformed financial economics from descriptive collection observations rigorous theoretical discipline Markowitz-Sherpa-CAPM framework underpins modern portfolio theory Sharpe s reward-to-risk ratio performance benchmarks MM theorem fundamental building block corporate finance Capital structure trade-off theories extend MM incorporating taxes bankruptcy costs agency problems Finance textbooks organize theory around these three pillars their work established intellectual coherence financial markets theory Corporate governance compensation stock options derivative hedging all rest on foundations laid by these three laureates

\section{Modern Perspectives Today}
Markowitz mean-variance optimization foundation Modern Portfolio Theory MPT expanded incorporating higher moments alternative risk measures downside risk Value-at-Risk Expected Shortfall Sharpe alpha-beta decomposition central performance evaluation mutual funds hedge funds pension plans Beta factors underpin factor investing style rotating strategies Miller MM theorem remains core corporate finance teaching Debt-equity leverage decisions considered along side profitability investment decisions CAPM modified incorporating multiple factors Fama's-French three-factor model four-factor five-factor models all modify original Sharpe framework capturing additional risk dimensions derivatives pricing Black-Scholes model builds upon Sharpe-Market concept of risk-neutral valuation Corporate finance theory continues exploring effects financing decisions on investment governance MM theorem provides benchmark against which imperfections measured Real-world deviations informed regulation corporate governance codes Investment advisors apply Markowitz diversification principles daily portfolio construction Risk managers implement VaR measures based on Sharpe insights Corporate treasurers make financing decisions informed by MM theorem extensions Each year thousands new papers extending Markowitz Sharpe Miller frameworks their work remains vibrant research frontier Financial education centers around their contributions every finance student learns their names as founders of modern financial economics
